\scanpage{144}
% archive foliation: 143

\uline{Construction abstraite d'une théorie de motifs semi-simples.}

On part d'une catégorie \(\mathcal V\)
\deleted{\unsure{pré}}``additive'', \deleted{\unsure{i.e.}}
\inserted{en particulier} les
\(\deleted{\operatorname{Hom}}\inserted{\operatorname{Fl}}(X,Y)\)
\deleted{\unsure{sont}} munis de une structure de
groupe abélien, les accouplements de composition
étant bilinéaires.
\deleted{[On ne suppose pas qu'il
existe un objet nul, et des sommes \unsure{pour deux}
\unsure{ou deux|de deux} objets......]}.
On s'intéresse, pour
toute catégorie \uline{abélienne} \(\mathcal C\), aux foncteurs
\deleted{contravariants} ``additifs''
(i.e. commutant aux sommes
finies, --- \inserted{\unsure{ce|ou} qui revient au \(=\)} aux pr. directs finis)
\[
H\deleted{\unsure{\,\mathcal V}}:
\mathcal V^{\unsure{\circ|*}}\longrightarrow\mathcal C
\]
et on cherche un foncteur universel
de cette nature
\[
\underline h:\mathcal V^{\unsure{\circ|*}}\longrightarrow\mathcal M(\mathcal V)=\mathcal M
\]
foncteur additif dans catégorie abélienne \(\mathcal M\)) tel
que pour toute catégorie abélienne \(\mathcal C\),
l'application
\[
\begin{gathered}
T\longmapsto T\circ\underline h\\
\operatorname{Fonc\,exacts}(\mathcal M,\mathcal C)
\longrightarrow\operatorname{Fonc\,add}(\mathcal V^{\unsure{\circ|*}},\mathcal C)
\end{gathered}
\]
soit bijective (ce qui caractérise \(\mathcal M\) à \deleted{isomorphisme}
isom près), ou le cas échéant, que
\[
\underline{\operatorname{Fonc\,exacts}}(\mathcal M,\mathcal C)
\mathbin{-}
\underline{\operatorname{Fonc\,add}}(\mathcal V^{\unsure{\circ|*}},\mathcal C)
\]
soit une équivalence de catégories, \unsure{ce qui|c à qui}
